
\subsubsection{Key Findings}

\textbf{Finding 1: Epistemic uncertainty threshold exists at 8B scale.} MC dropout $k\geq 3$ required for AUROC $\geq$ 0.70, while zero-cost methods (temperature scaling, conformal prediction) fall below threshold. This suggests post-hoc calibration (aleatoric uncertainty) is insufficient when base model has complex miscalibration patterns. Epistemic uncertainty---captured via stochastic forward passes---provides stronger signal for rejecting incorrect predictions. The mechanism matters at 8B scale: temperature scaling rescales logits assuming monotonic confidence-correctness relationship, but TruthfulQA adversarial design breaks this assumption. MC dropout samples from approximate Bayesian posterior, directly quantifying regions where model is uncertain.

\textbf{Finding 2: MC dropout k=3 is an efficiency point.} k=3 achieves 0.704 AUROC at $3\times$ cost. This contrasts with computer vision conventions (k=30) and suggests Bayesian approximation converges faster at 8B model scale. For practitioners, k=3 should be considered for applications requiring $\geq 0.70$ AUROC. The marginal AUROC gain from k=5 (+0.008) rarely justifies doubling cost from $3\times$ to $5\times$, unless application is high-stakes.

\textbf{Finding 3: Pareto frontier validates budget-aware UQ selection.} Five methods occupy distinct positions. No universal dominance across budget constraints. This shifts evaluation from ''which method is best'' to ''which method for my budget.'' Zero-cost methods are acceptable for applications with <$3\times$ budget constraint AND tolerance for <0.70 AUROC.

\subsubsection{Limitations}

\textbf{Limitation 1: 8B model scale only.} Results are specific to Llama-3.1-8B-Instruct. Larger models (70B, 405B) may have better base calibration, potentially enabling zero-cost methods to exceed 0.70 threshold. Smaller models (GPT-2, 1B) may require higher k for threshold crossing. 8B scale is realistic for production deployments (latency/memory constraints). Our findings establish feasibility at this scale and identify k=3 efficiency point. Future work: repeat experiments on Llama-3.1-70B-Instruct to test if zero-cost methods become competitive at larger scale.

\textbf{Limitation 2: Single benchmark (TruthfulQA).} Task-dependent AUROC thresholds exist. MMLU may be easier, GSM8K harder. The 0.70 threshold may not generalize across domains. TruthfulQA adversarial design provides challenging testbed (31\% base accuracy). Methods working here likely transfer to easier tasks. Threshold can be task-adjusted. Future work: cross-benchmark validation (MMLU, GSM8K, HaluEval) to characterize task-dependent thresholds and method ranking stability.

\textbf{Limitation 3: Limited statistical power.} Single seed for full-scale validation. MC k=10 versus k=5 difference (+0.006 AUROC) not assessed for statistical significance. Small differences suggest result is genuine though underpowered. Future work: multiple seeds for statistical power in pairwise comparisons.

\textbf{Limitation 4: Cross-dataset calibration transfer gap.} Conformal prediction AUROC 0.695 (-0.005 below threshold) may reflect domain shift. In-distribution calibration expected to improve AUROC. Coverage guarantees still hold per conformal theory. AUROC is quality metric, not validity metric. Future work: use TruthfulQA 40\% calibration split instead of cross-dataset calibration to close gap.

\subsubsection{Broader Impact}

Budget-aware UQ selection enables more practitioners to deploy selective prediction in production. k=3 efficiency point reduces barrier to entry for resource-constrained teams. Practitioners may over-optimize for cost at expense of safety. Applications requiring high confidence (medical, legal) should not use k=3 if k=5 is affordable. Our framework guides trade-offs but does not mandate specific choices. Clear communication that 0.70 threshold is baseline for viable selective prediction, not guarantee of safety. High-stakes applications should use $k\geq 5$ and validate on task-specific benchmarks.
